\DocumentMetadata{
  pdfversion=2.0,
  pdfstandard=ua-2,
  lang=en-US,
  tagging=on,
  tagging-setup={math/setup=mathml-SE,tabsorder=structure}
}
\documentclass[11pt,letterpaper]{article}
\usepackage[margin=0.68in,includefoot]{geometry}
\usepackage{newcomputermodern}
\usepackage{amsmath,amscd,mathtools}
\usepackage{tikz,adjustbox}
\providecommand{\square}{\mdlgwhtsquare}
\usepackage[hidelinks]{hyperref}
\hypersetup{
  pdftitle={Complex Analysis PhD Exam, January 2020},
  pdfauthor={University of Florida Department of Mathematics},
  pdfsubject={Graduate mathematics examination},
  pdfdisplaydoctitle=true,
  bookmarksopen=true
}
\setcounter{secnumdepth}{0}
\setlength{\parindent}{0pt}
\setlength{\parskip}{0.28em}
\setlength{\emergencystretch}{3em}
\setlength{\leftmargini}{2.15em}
\setlength{\leftmarginii}{2.65em}
\setlength{\leftmarginiii}{3em}
\renewcommand{\labelenumi}{\textbf{\theenumi.}}
\pagestyle{empty}
\raggedbottom
\allowdisplaybreaks
\newenvironment{examproblems}[1]{%
  \begin{enumerate}%
  \setcounter{enumi}{#1}%
  \setlength{\topsep}{0.35em}%
  \setlength{\partopsep}{0pt}%
  \setlength{\itemsep}{0.48em}%
  \setlength{\parsep}{0.18em}%
}{\end{enumerate}}
\newenvironment{examsubparts}{%
  \begin{enumerate}%
  \setlength{\topsep}{0.18em}%
  \setlength{\partopsep}{0pt}%
  \setlength{\itemsep}{0.16em}%
  \setlength{\parsep}{0.1em}%
}{\end{enumerate}}
\newenvironment{examinstructions}{%
  \par\begingroup\itshape
}{%
  \par\endgroup\vspace{0.18em}
}
\makeatletter
\renewcommand\section{\@startsection{section}{1}{0pt}%
  {-1.25ex plus -.25ex minus -.1ex}%
  {0.55ex plus .1ex}%
  {\normalfont\large\bfseries}}
\renewcommand{\@maketitle}{%
  \newpage\null\vskip -1em%
  {\footnotesize\bfseries
    \href{https://www.ufl.edu/}{University of Florida}, \href{https://math.ufl.edu/}{Department of Mathematics}\par}%
  \vskip -0.5em%
  \tagstructbegin{tag=Title}%
    {\LARGE\bfseries\href{https://gma.math.ufl.edu/exams/complex-analysis-phd/}{Complex Analysis PhD Exam}, January 2020\par}%
  \tagstructend%
  \vskip 0.25em%
  \tagmcbegin{artifact}%
    \hrule height 1pt%
  \tagmcend%
  \vskip 1em%
}
\makeatother
\title{Complex Analysis PhD Exam, January 2020}
\author{}
\date{}
\begin{document}
\maketitle
\thispagestyle{empty}
\begin{examinstructions}
Answer SIX questions in detail. State carefully any results used without proof. Throughout, \(\mathbb{D}\) is the open unit disc in~\(\mathbb{C}\).
\end{examinstructions}
\begin{examproblems}{0}
\item
When \(n>1\) is a positive integer, evaluate the integral 
\[
\int_0^\infty \frac{1}{1+x^n}\,dx.
\]
\item
True or false? There exists a holomorphic function \(f:\mathbb{D}\to\mathbb{C}\) such that \(f(1/n)=f(-1/n)=1/n^3\) for each integer \(n>1\). Prove or disprove.
\item
The functions~\(f\) and~\(g\) are holomorphic and never zero in~\(\mathbb{D}\). Assume that 
\[
\frac{f'}{f}(2^{-n})=\frac{g'}{g}(2^{-n})
\]
 for each positive integer~\(n\). Deduce as much as possible about the relationship between~\(f\) and~\(g\).
\item
Let~\(f\) be holomorphic in \(\mathbb{D}\setminus\{0\}\) and assume that its derivative \(f'\) extends holomorphically to~\(\mathbb{D}\). Must~\(f\) itself extend holomorphically to~\(\mathbb{D}\)? Explain.
\item
Let the continuous function \(f:\mathbb{D}\to\mathbb{C}\) be holomorphic except (perhaps) on the interval \([-1/2,1/2]\). Prove that~\(f\) is in fact holomorphic on~\(\mathbb{D}\).
\item
Let the entire function~\(f\) never take purely imaginary values. Without using a theorem of Picard, prove that~\(f\) is constant.
\item
Let the entire function~\(f\) satisfy \(|f(z)|\leq A+B|z|^M\) whenever \(|z|\geq N\), where \(A,B,M,N\) are positive constants. Prove that~\(f\) is a polynomial.
\item
Let~\(f\) be holomorphic in~\(\mathbb{D}\). Assume that for each \(z\in\mathbb{D}\) there exists \(n=n_z\in\mathbb{N}\) such that \(f^{(n)}(z)=0\). Prove that~\(f\) is a polynomial.
\item
Prove that if~\(z\) and~\(w\) lie in the (open) right half-plane then 
\[
\int_0^\infty e^{-wt}t^{z-1}\,dt=w^{-z}\Gamma(z)
\]
 by integrating \(e^{-\zeta}\zeta^{z-1}\) around a suitable closed contour.
\end{examproblems}
\end{document}
